\chapter{\textbf{Large strain and hyperelasticity}}
\label{chap:largestrain}

\section*{Outline}
\addcontentsline{toc}{section}{\texorpdfstring{\hskip14mm}{} Outline}

Everything so far assumed small strains: the equilibrium equations were
written on the undeformed geometry, and strain was the symmetric part of the
displacement gradient. A rubber block stretched to three times its length,
or a thin shell that buckles and then keeps deforming, satisfies neither
assumption.

This chapter covers what changes. It is in three parts: the strain measure,
which is a modelling choice \texttt{Cast3M} makes you declare; the
constitutive laws for rubber-like materials, which are built into the code;
and the \op{pasapas} options that drive the whole thing. It closes on contact
and post-buckling, where large strain earns its keep.

The code fragments are extracts rather than complete files: they carry over
objects from the chapters before (the cylinder of chapter~\ref{chap:meshing},
the buckling modes of chapter~\ref{chap:eigen}) and from the shipped examples
they are modelled on, whose names are given as each one appears.

The theory is only sketched here. For the continuum mechanics, Holzapfel is
the single best companion~\cite{Holzapfel:2000}; for the finite element side
-- linearisation, tangent operators, implementation -- Bonet and
Wood~\cite{BonetWood:2008}; and Ogden~\cite{Ogden:1984} for the theory of
nonlinear elasticity itself.

\section{What changes in large strain}
\label{sec:ls:changes}

Three things, and they are independent of each other:

\begin{itemize}
\cindex{large strain}\cindex{large displacement}%
\cindex{finite rotation}\cindex{objectivity}%
\item \textbf{the geometry moves.} Equilibrium must be written on the
  deformed configuration, which has to be updated as the computation
  proceeds. This is the \emph{large displacement} part, and it is what
  \texttt{'GRANDS\_DEPLACEMENTS'} switches on;
\item \textbf{the strain measure.} With finite rotations, the linearised
  strain is no longer objective -- a rigid body rotation produces fictitious
  strain. A different measure is needed, and \texttt{Cast3M} lets you choose
  which;
\item \textbf{the constitutive law.} Hooke's law relates a small stress to a
  small strain. A rubber needs a strain-energy function of the deformation
  invariants instead.
\end{itemize}

\noindent You can have the first without the others -- a slender beam in
large deflection but small strain is the classic case -- which is why they
are separate switches.

\subsection{Choosing the strain measure}

\begin{castemcom}
\sindex{opti 'EPSI'}\sindex{mode \ldots 'EPSI'}%
\cindex{strain measure}\cindex{Green--Lagrange strain}%
\cindex{Jaumann derivative}\cindex{Truesdell derivative}%
The strain measure is selected by the keyword \texttt{'EPSI'}, either on the
model with \op{mode} or globally with \op{opti}. The five values are
\texttt{'LINEAIRE'}, \texttt{'QUADRATIQUE'} (the default),
\texttt{'TRUESDELL'}, \texttt{'JAUMANN'} and \texttt{'UTILISATEUR'}.

In every case the code hands the constitutive law the gradient of the
displacement with respect to the initial state at the beginning and at the
end of the step; the keyword decides what is built from it. With
\texttt{'UTILISATEUR'} the strains passed to the law are linear but computed
on the half-step configuration, which is the usual choice for the
hyperelastic laws below.
\end{castemcom}
\begin{castemlines}
\begin{verbatim}
* globally ...
opti epsi 'QUADRATIQUE';

* ... or per model
mod = vol mode mecanique
          elastique isotrope
          'EPSI' 'UTILISATEUR';
\end{verbatim}
\end{castemlines}

\begin{castemcom}
\sindex{'GRANDES\_DEFORMATIONS' (obsolete)}%
\sindex{'GRANDES\_ROTATIONS' (obsolete)}%
\textbf{A warning for anyone with older files.} The \op{pasapas} table
indices \texttt{'GRANDES\_DEFORMATIONS'} and \texttt{'GRANDES\_ROTATIONS'}
are obsolete. They are not ignored: current versions detect them and
\emph{stop the computation} with an explicit message telling you to use the
\texttt{'EPSI'} keyword of \op{opti} or \op{mode} instead.

If you inherit a file from a colleague that aborts immediately with a message
about \emph{grandes rotations}, this is why.

\medskip
\textbf{This affects the published material too.} The Club Cast3M 2011
presentation~\cite{Gornet:club2011}, which is otherwise the best guide to the
hyperelastic laws, drives its examples with
\texttt{'GRANDES\_DEFORMATIONS' = VRAI}. That line was correct when the
slides were written and stops the computation today. Read the slides for the
method and take the syntax from the shipped examples, which are kept up to
date.
\end{castemcom}
\begin{castemlines}
\begin{verbatim}
* -- obsolete, aborts the run --
* tab1 . 'GRANDES_DEFORMATIONS' = vrai;
* tab1 . 'GRANDES_ROTATIONS'    = vrai;

* -- current --
tab1 . 'GRANDS_DEPLACEMENTS' = vrai;
tab1 . 'HYPOTHESE_DEFORMATIONS' =
       mot 'UTILISATEUR';
\end{verbatim}
\end{castemlines}

\section{Hyperelastic constitutive laws}
\label{sec:ls:hyper}

\cindex{hyperelasticity}\cindex{strain energy density}%
\cindex{rubber, constitutive laws}\cindex{invariants of the strain}%
A hyperelastic material is defined by a strain energy density $W$, from which
the stress derives. For an incompressible isotropic material $W$ is a
function of the first two invariants $I_{1}$, $I_{2}$ of the left
Cauchy--Green tensor, and the many models in the literature differ in the
form chosen for that function.

\subsection*{Which models are in \texttt{Cast3M}, and who put them there}

\cindex{Gornet, L.}\cindex{Marckmann, G.}\cindex{Verron, E.}%
\cindex{Club Cast3M}\cindex{Mullins effect}%
The hyperelastic laws in \texttt{Cast3M} are the contribution of
\textbf{Laurent Gornet} (Centrale Nantes, Institut GeM). This is not
folklore: he is credited by name in the notice of the \op{mate} operator, and
the source of the Gornet--Desmorat routine carries the header

\begin{center}\small\ttfamily
MODELE HYPERELASTIQUE GORNET-DESMORAT\\
Contribution de Laurent Gornet -- Ecole Centrale de Nantes (2009)
\end{center}

\noindent followed by a pointer to his Club Cast3M presentation. He has
returned to the Club five times between 2006 and 2017, and the four
presentations that matter here are, in order of usefulness:

\begin{itemize}
\item \textbf{Club Cast3M 2011}~\cite{Gornet:club2011} -- \emph{the
  implementation reference}. The only one that descends to operator level:
  how to declare the GD and GDM laws through \op{mode} with
  \texttt{'NUME\_LOI' 33}, how to pass the parameters with \op{mate}, and how
  to drive the whole thing with \op{pasapas}, for plane stress, plane strain,
  axisymmetric and 3D. Read this one before writing your own file;
\item \textbf{Club Cast3M 2009}~\cite{Gornet:club2009} -- \emph{which law to
  pick}. A systematic comparison from Mooney--Rivlin (2 constants) through
  Gent--Thomas, Hart-Smith, Biderman and Arruda--Boyce to Ogden (6 constants)
  and GD (3 constants), validated against the Treloar (1944) and Kawabata
  (1981) data in simple extension, equibiaxial, pure shear and general
  biaxial loading;
\item \textbf{Club Cast3M 2006}~\cite{Gornet:club2006} -- \emph{how a user law
  gets into the code at all}: the \texttt{UMAT} methodology, and the bridge
  from a \textsc{fortran} material routine into the \textsc{gibiane} and
  \textsc{esope} environments. The place to start if you ever need to add a
  constitutive law of your own;
\item \textbf{Club Cast3M 2017}~\cite{Gornet:club2017} -- the \texttt{@Torsion}
  procedure for beam cross-sections, discussed in
  section~\ref{sec:elas:torsion}.
\end{itemize}

\noindent The model itself is published as
Gornet, Marckmann, Desmorat and Charrier~\cite{Gornet:2012}, and its damage
lineage -- the network alteration theory behind the Mullins effect -- as
Marckmann \emph{et al.}~\cite{Marckmann:2002}.

For the choice between models the standard reference is the systematic
comparison by \textbf{Marckmann and Verron}~\cite{Marckmann:2006}, also from
the Nantes group, which ranks twenty hyperelastic models against the classical
Treloar and Kawabata data sets. If you need one paper to decide which law to
use for a given rubber and a given strain range, that is the one.

\medskip
\noindent Two practical consequences of how these laws are implemented:

\begin{itemize}
\item \textbf{the base GD law ships with the code, the Mullins extension does
  not.} \texttt{'NUME\_LOI' 33} gives you GD; GDM, which adds isotropic and
  directional Mullins damage, is described in the 2011 presentation but is not
  in the distribution, so using it means implementing it yourself by the
  \texttt{UMAT} route of the 2006 presentation;
\item \textbf{in large strain you cannot mix law families.} The source of the
  GD routine states it plainly: in large deformation under \op{pasapas} the
  model may contain \texttt{UMAT}-type laws only -- objective derivatives
  cannot be mixed within one model. If a large-strain computation combines a
  hyperelastic part with a law from elsewhere in the catalogue, this is why it
  will not run.
\end{itemize}

\cindex{Mooney--Rivlin model}\cindex{neo-Hookean model}%
\cindex{Gornet--Desmorat model}\cindex{Hart--Smith model}%
\cindex{Biderman model}\cindex{Arruda--Boyce model}\cindex{8-chain model}%
\cindex{Ogden model (absent from Cast3M)}%
\begin{optable}{Hyperelastic models available}
31 & Mooney--Rivlin (and neo-Hooke as the special case $C_2=0$) \\
32 & neo-Hookean, compressible \\
33 & GD -- Gornet--Desmorat \\
34 & Hart--Smith \\
35 & Biderman \\
36 & 8-chain (Arruda--Boyce) \\
\end{optable}

\noindent There is \textbf{no Ogden and no Signorini model} in
\texttt{Cast3M}. The notice positions GD as the substitute: it is expressed
from the invariants and is stated to perform comparably to a five-parameter
Ogden model.

\subsection*{Declaring a hyperelastic model}

\sindex{'NUME\_LOI'}\sindex{'C\_MATERIAU'}\sindex{'NON\_LINEAIRE' 'UTILISATEUR'}%
\cindex{user material law}\cindex{UMAT interface}%
This is the part whose syntax is unlike anything else in the booklet. The
models are not separate keywords of \op{mode}; they are \emph{user law
numbers} passed through the \texttt{'NON\_LINEAIRE' 'UTILISATEUR'} mechanism.

\begin{castemcom}
\op{mots} builds the list of parameter names the law expects, which is passed
as \texttt{'C\_MATERIAU'}; \texttt{'NUME\_LOI'} selects the model from the
table above.

\medskip
\textbf{Note the padding.} Component names are four characters, and the
shipped examples write them padded: \texttt{'NU~~'} with two trailing
spaces, \texttt{'C1~~'}, \texttt{'C2~~'}. \op{mots} pads them for you, so
what actually matters is that the list given to \texttt{'C\_MATERIAU'} and
the names given to \op{mate} agree with each other; a mismatch between the
two is the usual reason the law reports an unknown parameter.

\textbf{And declare \texttt{'YOUN'} and \texttt{'NU'} even though the model
does not use them}: they are needed by the convergence test inside
\op{pasapas}. For an incompressible model take $\nu$ close to $0.5$; the
shear modulus is then $\mu = E/\bigl(2(1+\nu)\bigr) = 2(C_{1}+C_{2})$, which
fixes $E$.
That value is the initial modulus, and it is a lower bound -- at large
stretch the tangent modulus rises, and the value may have to be increased to
get convergence.
\end{castemcom}
\begin{castemlines}
\begin{verbatim}
* Mooney-Rivlin, incompressible,
* plane stress
lcmat = mots 'YOUN' 'NU  '
             'C1  ' 'C2  ';

mod = su mode 'MECANIQUE'
              'ELASTIQUE' 'ISOTROPE'
              'NON_LINEAIRE' 'UTILISATEUR'
              'NUME_LOI' 31
              'C_MATERIAU' lcmat;

* Treloar's vulcanised rubber, in Pa
c1 = 0.183e6;  c2 = 0.0034e6;
xnu = 0.4995;
you = 2. * (1. + xnu) * 2. * (c1 + c2);

mat = mod mate 'YOUN' you 'NU  ' xnu
               'C1  ' c1  'C2  ' c2;
\end{verbatim}
\end{castemlines}

\cindex{incompressibility}\cindex{penalty coefficient}%
\noindent For a quasi-incompressible formulation an extra penalty
coefficient \texttt{'D'} appears in the parameter list, governing the
hydrostatic part $W(J)=\frac{5}{D}(J-1)^{2}$ of the energy; its default is
$10^{-4}$. The Gornet--Desmorat law takes \texttt{'H1'}, \texttt{'H2'},
\texttt{'H3'}; Hart--Smith takes \texttt{'G'}, \texttt{'K1'}, \texttt{'K2'};
the 8-chain model takes \texttt{'NKT'} and \texttt{'VN~~'}.

\section{Driving the computation with \texttt{pasapas}}

The table is the one from chapter~\ref{chap:plasticity}, with a handful of
extra entries.

\textbf{A version caveat before the table.} \texttt{'HYPOTHESE\_DEFORMATIONS'}
and the \texttt{'MI\_PAS'}/\texttt{'FIN\_PAS'} values of
\texttt{'LAGRANGIEN'} are recent additions: they are in the notice of current
versions but not in distributions of a few years ago, where the strain measure
could only be set through \op{opti 'EPSI'} or \op{mode \ldots 'EPSI'} and
\texttt{'LAGRANGIEN'} took only \texttt{'TOTAL'} or \texttt{'REACTUALISE'}.
Setting the strain measure on the model, as in section~\ref{sec:ls:changes},
works on every version and is the portable choice.

\sindex{'GRANDS\_DEPLACEMENTS'}\sindex{'HYPOTHESE\_DEFORMATIONS'}%
\sindex{'REAC\_GRANDS'}\sindex{'K\_SIGMA'}\sindex{'PREDICTEUR'}%
\sindex{'LAGRANGIEN'}\sindex{'CONVERGENCE\_FORCEE'}%
\begin{optable}{\texttt{pasapas} entries for large strain}
'GRANDS\_DEPLACEMENTS' & VRAI to update the geometry at each iteration.
  Also sets \texttt{'K\_SIGMA'} to true \\
'HYPOTHESE\_DEFORMATIONS' & the strain measure, as in \texttt{'EPSI'} above \\
'REAC\_GRANDS' & strain increment beyond which the stiffness matrix is
  reassembled (default $10^{-1}$) \\
'K\_SIGMA' & VRAI to add the geometric stiffness to the iteration operator \\
'PREDICTEUR' & \texttt{'HPP'} initialises the step with a small-strain
  computation \\
'LAGRANGIEN' & the configuration the constitutive law is applied on:
  \texttt{'MI\_PAS'} (also spelled \texttt{'REACTUALISE'}, the default),
  \texttt{'FIN\_PAS'}, or \texttt{'TOTAL'} for the initial geometry \\
\end{optable}

\begin{castemcom}
A complete driver, following the shipped example
\texttt{Mooney\_LRGTreloar\_Cisaillementsimple}: a simple shear of a
Mooney--Rivlin block in plane stress, which has an analytical solution to
check against.

\op{'REAC\_GRANDS'} is the one to tune if the computation struggles: it sets
how much strain may accumulate before the stiffness is rebuilt. Smaller means
more reassembly and more cost, but better convergence.

\medskip
The other lever is the time step. Unlike the elastoplastic computation of
chapter~\ref{chap:plasticity}, where the step size mainly controls accuracy,
here it controls whether the thing converges at all -- a hyperelastic model
at large stretch has a strongly varying tangent modulus.
\end{castemcom}
\begin{castemlines}
\begin{verbatim}
tab1 = table;
tab1 . 'MODELE' = mod;
tab1 . 'CARACTERISTIQUES' = mat;
tab1 . 'BLOCAGES_MECANIQUES' = bltot;
tab1 . 'CHARGEMENT' = chartot;

tab1 . 'GRANDS_DEPLACEMENTS' = vrai;
tab1 . 'HYPOTHESE_DEFORMATIONS' =
       mot 'UTILISATEUR';
tab1 . 'REAC_GRANDS' = 0.4;
tab1 . 'CONVERGENCE_FORCEE' = faux;

tab1 . 'TEMPS_CALCULES' =
       prog t_deb pas 0.5 t_fin;
tab1 . 'TEMPS_SAUVES'   =
       prog t_deb pas 0.5 t_fin;

pasapas tab1;
\end{verbatim}
\end{castemlines}

\noindent \texttt{'CONVERGENCE\_FORCEE' = faux} is worth setting explicitly:
it makes \op{pasapas} stop rather than carry on from a step that did not
converge. In a large-strain computation a forced convergence does not give a
slightly wrong answer, it gives a meaningless one.

\subsection*{Example files to start from}

The distribution ships about twenty-five hyperelastic test cases, named
\texttt{<model><loading><hypothesis>} -- \texttt{trac} for traction,
\texttt{cis} for shear, \texttt{dp} for plane strain, \texttt{3d} for three
dimensions:

\begin{center}
\begin{tabular}{@{}>{\ttfamily}l@{\hspace{6mm}}p{85mm}@{}}
Mooney\_LRGTreloar\_Traction & Mooney--Rivlin, uniaxial tension, checked
  against the analytical solution \\
Mooney\_LRGTreloar\_Cisaillementsimple & the same law in simple shear;
  \textbf{the best single teaching example}, since it uses
  \texttt{'GRANDS\_DEPLACEMENTS'}, \texttt{'HYPOTHESE\_DEFORMATIONS'} and
  \texttt{'REAC\_GRANDS'} together \\
NeoHookeen\_Traction3D & compressible neo-Hooke, 3D \\
gdtract & Gornet--Desmorat, plane stress \\
harttrac & Hart--Smith, plane stress \\
\end{tabular}
\end{center}

\section{Contact}
\label{sec:ls:contact}

Large displacement and contact go together: as soon as bodies move far
enough to touch, the geometry on which contact is detected is the deformed
one, which is exactly what \texttt{'GRANDS\_DEPLACEMENTS'} maintains.

\begin{castemcom}
\dindex{impo}{contact condition}\sindex{impo 'MAIL'}%
\dindex{orie[nter]}{orient a surface}%
\cindex{contact}\cindex{unilateral contact}\cindex{friction, Coulomb}%
\cindex{master--slave contact}%
Contact is set up in two steps. \op{impo 'MAIL'} builds an intermediate mesh
of contact elements between two lines (2D, \texttt{SEG2}) or two surfaces
(3D, \texttt{TRI3}); \op{mode \ldots 'CONTACT' 'UNILATERAL'} turns that mesh
into a model, which is then concatenated to the mechanical model and handed
to \op{pasapas} in the usual way.

\medskip
\textbf{Orientation is everything.} The two surfaces must be oriented facing
each other with \op{orie}, and they must be \texttt{TRI3} in 3D -- hence the
\op{chan 'TRI3'}. Getting the normals wrong is the single most common cause
of a contact computation that does nothing, or that locks solid.

Friction is available through \op{mode \ldots 'CONTACT' 'FROTTANT'
'COULOMB'}.
\end{castemcom}
\begin{castemlines}
\begin{verbatim}
* orient the two surfaces so that
* they face each other
sol2 = orie (chan 'TRI3' sol1)
            (0. 0. 1.);
scc1 = orie (chan 'TRI3' sc1)
            (0. 0. -1.);

mcont = impo 'MAIL' sol2 scc1;
modc  = mode mcont 'CONTACT'
             'UNILATERAL';

tab1 . 'MODELE' = mod et modc;
\end{verbatim}
\end{castemlines}

\subsection*{Self-contact}

A body that folds onto itself -- a seal being crushed, a thread under
torsion, a wrinkling membrane -- contacts \emph{itself}, so the master and
slave surfaces are parts of the same mesh. \texttt{Cast3M} handles this
through a second, optional criterion on \op{impo 'MAIL'}:

\begin{castemcom}
\cindex{self-contact (autocontact)}\sindex{crit2, self-contact criterion}%
\op{impo 'MAIL'} accepts up to two numerical criteria. The first,
\op{crit}, limits the generated contact elements to node pairs closer than
that distance. The second, \op{crit2}, is available in 3D and is the minimum
distance between an element and the candidate nodes: it exists specifically
to stop an element from coming into contact with itself. That is the
self-contact mechanism.

\medskip
A warning: setting \op{crit} too large generates an enormous contact mesh and
the computation slows to a crawl; too small and contact is simply not
detected. It is the parameter to tune first when a self-contact computation
misbehaves.
\end{castemcom}
\begin{castemlines}
\begin{verbatim}
* self-contact: the same surface
* given as both arguments, with the
* second criterion excluding an
* element contacting itself
surf = orie (chan 'TRI3' peau)
            (0. 0. 1.);

mauto = impo 'MAIL' crit crit2
             surf surf;
modc  = mode mauto 'CONTACT'
             'UNILATERAL';
\end{verbatim}
\end{castemlines}

\noindent \textbf{On the torsion example.} There is no self-contact example
with a thread under torsion in the distribution. I searched the shipped test
cases, the official test-case classification and the notices: the contact
cases include \texttt{Contact2D}, \texttt{Contact2Djeu},
\texttt{Contact2Djeufaible}, \texttt{Contact3D}, \texttt{Contact3Djeu},
\texttt{Coulomb3D}, \texttt{contact2D-adhe} and \texttt{dynacontact}, and
none of them is self-contact or involves torsion. The only occurrence of the word \emph{autocontact} anywhere
is the \op{impo} notice quoted above.

What does exist, and may be the source of the recollection, is a
\texttt{@TORSION} procedure for the Saint-Venant torsion and warping of beam
cross-sections, presented at the Club Cast3M in 2017 -- by Laurent Gornet
again, the same author as the hyperelastic laws. It is a beam
cross-section tool and involves no contact.

So the capability is there and the ingredients are documented, but the worked
example would have to be built. The exercises below suggest how.

\section{Post-buckling: beyond the eigenvalue}

Chapter~\ref{chap:eigen} computed the load at which a perfect structure
bifurcates. It said nothing about what happens next, and for the cylinder
that is the interesting part: the real structure collapses well below the
eigenvalue, and the post-buckling path turns back on itself.

\cindex{post-buckling}\cindex{snap-through}%
\cindex{geometric imperfection}%
Following that path is a large-displacement computation, and it needs two
things the computations so far have not: a geometric imperfection, since a
perfect structure has no reason to leave the fundamental path, and an
algorithm able to follow a load--displacement curve that is not monotonic.
The shipped examples \texttt{gdep1}, \texttt{gdep3}, \texttt{continu\_gdep1}
and \texttt{continu\_snap} cover this, the last being the classical
snap-through problem.

\begin{castemcom}
The imperfection is introduced by adding a fraction of the first buckling
mode to the mesh -- exactly the \op{plus} construction of
section~\ref{sec:mesh:deform}, with the mode shape of
chapter~\ref{chap:eigen} as the displacement field.

\op{maxi} with the keyword \texttt{'ABS'} returns the largest value of the
field in absolute value over all its components, which is what you want to
normalise a mode shape with: for the axially compressed cylinder the buckling
mode is mostly radial, so normalising on a single component such as
\texttt{'UZ'} would scale it wrongly.

A tenth of the thickness is a realistic amplitude for a manufactured shell.
The computed collapse load is then a function of that amplitude, and the
curve of one against the other is the imperfection sensitivity of the
structure.
\end{castemcom}
\begin{castemlines}
\begin{verbatim}
* first buckling mode, from the
* buckling computation of the
* eigenvalue chapter
mode1 = stab . 1 . 'DEPL';

* scale it to a fraction of the
* thickness and perturb the mesh
ampl  = 0.1 * 0.002;
nrm   = maxi mode1 'ABS';
imper = (ampl / nrm) * mode1;

cylimp = cyl plus imper;
\end{verbatim}
\end{castemlines}

\section*{Summary}
\addcontentsline{toc}{section}{\texorpdfstring{\hskip14mm}{} Summary}

Large strain adds three independent choices to the computations of the
previous chapters, and \texttt{Cast3M} keeps them separate. The geometry is
updated by \texttt{'GRANDS\_DEPLACEMENTS'}; the strain measure is chosen by
\texttt{'EPSI'} on \op{opti} or \op{mode}, \emph{not} by the obsolete
\texttt{'GRANDES\_DEFORMATIONS'} index, which now aborts the run; and the
constitutive law is selected through \texttt{'NON\_LINEAIRE' 'UTILISATEUR'}
with a \texttt{'NUME\_LOI'} number, with six hyperelastic models available.

The driver remains the \op{pasapas} table. Contact enters it as an additional
\emph{model}, built from \op{impo 'MAIL'} on two oriented surfaces, and
self-contact is the same construction with one surface given twice and the
second criterion \op{crit2} excluding self-pairs.

For choosing a rubber model, Marckmann and Verron~\cite{Marckmann:2006}
compare twenty of them against the classical data; the laws in
\texttt{Cast3M} itself are due to Gornet, whose Club Cast3M 2011
presentation~\cite{Gornet:club2011} is the implementation reference.

\section*{Exercises}
\addcontentsline{toc}{section}{\texorpdfstring{\hskip14mm}{} Exercises}

\begin{enumerate}
\item \textbf{Uniaxial tension of a rubber.} Run
  \texttt{Mooney\_LRGTreloar\_Traction} and reproduce the nominal
  stress--stretch curve up to $\lambda=7$. Compare with the analytical
  Mooney--Rivlin result
  $P = 2(\lambda - \lambda^{-2})(C_{1} + C_{2}\lambda^{-1})$. Where does the
  model start to deviate from Treloar's data, and which of the models of
  section~\ref{sec:ls:hyper} would do better there?

\item \textbf{Comparing models.} Fit Mooney--Rivlin, Hart--Smith and
  Gornet--Desmorat to the same uniaxial data and then use each to predict
  \emph{simple shear}. This is the test that separates them -- a model fitted
  on one loading mode need not predict another. Compare with the conclusions
  of~\cite{Marckmann:2006}.

\item \textbf{The incompressibility limit.} Rerun the Mooney--Rivlin tension
  with $\nu = 0.49$, $0.499$, $0.4999$ and $0.49999$. At what point does the
  computation stop converging, and what does that tell you about why a
  quasi-incompressible formulation with a penalty coefficient \texttt{'D'}
  exists?

\item \textbf{Hertz contact.} Build the contact of a cylinder on a plane in
  2D, starting from \texttt{Contact2D}, and compare the contact pressure
  distribution and the half-width of the contact patch with the analytical
  Hertz solution. Then halve the mesh size in the contact zone and see how
  much the peak pressure changes -- contact pressures converge slowly.

\item \textbf{Towards self-contact.} Build the simplest self-contact problem
  you can: a thin strip, clamped at one end, pushed at the other until it
  folds over and touches itself. Use \op{impo 'MAIL'} with the surface given
  twice, and tune \op{crit} and \op{crit2} until the fold is detected without
  the contact mesh exploding. This is the exercise that would turn
  section~\ref{sec:ls:contact} into a worked example -- and it is a reasonable
  small project rather than a half-hour exercise.

\item \textbf{Imperfection sensitivity of the cylinder.} Using the
  imperfection construction above, compute the collapse load of the cylinder
  for imperfection amplitudes of $0$, $0.1t$, $0.5t$ and $t$, and plot it
  against the eigenvalue of chapter~\ref{chap:eigen}. The fall is the reason
  shell design codes do not use the linear buckling load.
\end{enumerate}
